\chapter{Results}
\label{ch:results}

This chapter consolidates the quantitative results produced across
Chapters~\ref{ch:sim-stack}--\ref{ch:nav-integration} and evaluates
them directly against the target state defined in
Section~\ref{sec:problem-statement}: a Nav2 navigation stack that
reaches goal poses reliably on the physical QCar, with the
sim-to-real gap identified and, where practical, corrected. Results
are presented here without re-opening their diagnosis; interpretation
and generalisation are deferred to Chapter~\ref{ch:summary}.

\section{Actuator Calibration Accuracy}
\label{sec:results-throttle}

Table~\ref{tab:results-throttle-summary} restates the throttle
calibration validation of Section~\ref{sec:final-calibration}
alongside the uncalibrated baseline it replaced.

\begin{table}[htbp]
  \centering
  \caption{Commanded-versus-real cruise speed, before and after calibration.}
  \label{tab:results-throttle-summary}
  \begin{tabular}{@{}p{3.2cm}p{5.8cm}p{5.5cm}@{}}
    \toprule
    Commanded speed & Baseline ratio (uncalibrated $\text{THROTTLE\_GAIN}=1/3$) & Calibrated ratio (Eq.~\ref{eq:throttle-calibration}) \\
    \midrule
    \SI{0.3}{\meter\per\second} & $\approx$2.6--2.7 & 1.05 \\
    \SI{0.4}{\meter\per\second} & -- (not separately measured) & 1.05 \\
    \SI{0.5}{\meter\per\second} & -- (not separately measured) & 1.00 \\
    \bottomrule
  \end{tabular}
\end{table}

The commanded-to-real speed ratio was reduced from
$\approx$2.6--2.7$\times$ to within 5\% across the
\SIrange{0.3}{0.5}{\meter\per\second} range validated. Below
$\approx$\SI{0.25}{\meter\per\second}, both the calibrated and
uncalibrated mappings are unreliable in an open-loop sense, since the
vehicle's static-friction deadband (Section~\ref{sec:deadband-discovery})
sits in that range independently of gain calibration; commanding
\SI{0.2}{\meter\per\second} moved the vehicle in only some trials.

\section{Stop-Distance and Latency}
\label{sec:results-stop-distance}

Table~\ref{tab:results-stop-distance} restates the
\texttt{readMode=0} fix's effect (Section~\ref{sec:readmode-validation})
across all five tested conditions, alongside the residual physical
explanation for what remains.

\begin{table}[htbp]
  \centering
  \caption{Stop-latency and extra-distance results, before and after the \texttt{readMode=0} fix.}
  \label{tab:results-stop-distance}
  \begin{tabular}{@{}lccc@{}}
    \toprule
    Test condition & Stiction delay & Extra distance & Stop latency \\
    & (before $\to$ after) & (before $\to$ after) & (before $\to$ after) \\
    \midrule
    \SI{2}{\meter} @ \SI{0.3}{\meter\per\second} & \SIrange{2.09}{0.117}{\second} & \SIrange{0.520}{0.114}{\meter} & \SIrange{2.37}{0.598}{\second} \\
    \SI{2}{\meter} @ \SI{0.4}{\meter\per\second} & \SIrange{2.14}{0.001}{\second} & \SIrange{0.790}{0.114}{\meter} & \SIrange{2.40}{0.501}{\second} \\
    \SI{2}{\meter} @ \SI{0.5}{\meter\per\second} & \SIrange{2.14}{0.116}{\second} & \SIrange{0.923}{0.154}{\meter} & \SIrange{2.55}{0.550}{\second} \\
    \SI{1}{\meter} @ \SI{0.3}{\meter\per\second} & \SIrange{2.14}{0.120}{\second} & \SIrange{0.641}{0.118}{\meter} & \SIrange{2.44}{0.572}{\second} \\
    \SI{3}{\meter} @ \SI{0.3}{\meter\per\second} & \SIrange{2.08}{0.088}{\second} & \SIrange{0.583}{0.104}{\meter} & \SIrange{2.35}{0.509}{\second} \\
    \bottomrule
  \end{tabular}
\end{table}

Mean stiction delay across the five post-fix conditions is
$\approx$\SI{0.09}{\second} (down from $\approx$\SI{2.12}{\second}),
mean extra distance is $\approx$\SI{0.12}{\meter} (down from
$\approx$\SI{0.69}{\meter}), and mean stop latency is
$\approx$\SI{0.55}{\second} (down from $\approx$\SI{2.42}{\second}).
The post-fix residual is not further reducible with the current
open-loop throttle model: dividing extra distance by stop latency
recovers \SIrange{54}{61}{\percent} of cruise velocity in four of the
five conditions, consistent with an ordinary, roughly linear velocity
decay to zero over the measured latency window (a perfectly linear
decay predicts exactly 50\%), rather than indicating any further
unexplained delay mechanism.

\section{Sim-versus-Real Overshoot}
\label{sec:results-sim-vs-real}

Table~\ref{tab:results-sim-vs-real} restates the sim-versus-real
comparison of Section~\ref{sec:sim-vs-real}. Because this comparison
was run \emph{before} both the throttle calibration
(Section~\ref{sec:throttle-calibration}) and the buffering-defect fix
(Section~\ref{sec:readmode-bug}), it is reported here as the gap that
existed at that stage of the project, not as the final residual gap
between simulation and hardware -- Sections~\ref{sec:results-throttle}
and~\ref{sec:results-stop-distance} above are the corresponding
final, post-fix figures, but a matched six-point sim-versus-real
sweep was not re-run after both fixes were in place.
% TODO: if lab time allows, re-running the exact six-point sweep from
% Sec. 5.4 post-fix would let this results chapter report a genuine
% final sim-to-real gap number rather than only a pre-fix one.

\begin{table}[htbp]
  \centering
  \caption{Sim-versus-real stop-distance overshoot (measured before the fixes of Ch.~5).}
  \label{tab:results-sim-vs-real}
  \begin{tabular}{@{}lcc@{}}
    \toprule
    & Simulation & Real hardware \\
    \midrule
    Overshoot range (six-point sweep) & \SIrange{0.044}{0.067}{\meter} & \SIrange{0.58}{2.12}{\meter} \\
    Isolated \SI{2}{\meter}@\SI{0.3}{\meter\per\second} test & \SI{0.067}{\meter} extra & \SI{1.90}{\meter} extra \\
    Ratio (real / sim), isolated test & \multicolumn{2}{c}{$\approx$28$\times$} \\
    Speed sensitivity & weak & strong ($\propto v^3$, approximately) \\
    \bottomrule
  \end{tabular}
\end{table}

\section{Steering Response Accuracy}
\label{sec:results-steering}

Table~\ref{tab:results-steering-summary} restates the steering-gain
validation of Section~\ref{sec:steering-gain}, isolating the
multiplicative gain error from the zero-point trim already covered by
Section~\ref{sec:results-throttle}'s companion actuator.

\begin{table}[htbp]
  \centering
  \caption{Real wheel deflection versus kinematically commanded deflection, before and after \texttt{STEERING\_GAIN} correction.}
  \label{tab:results-steering-summary}
  \begin{tabular}{@{}lcc@{}}
    \toprule
    & Uncorrected (gain$=1$) & Corrected (\texttt{STEERING\_GAIN}$=1.39$) \\
    \midrule
    Residual error vs.\ kinematic prediction & $\approx$39\% & $\approx$5.6\% \\
    \bottomrule
  \end{tabular}
\end{table}

As with the throttle mapping, this is a per-unit, empirically fitted
correction rather than a property derivable from the vehicle's nominal
kinematic model alone -- a second, independent confirmation of this
thesis's broader finding that the QCar's actual actuator response
cannot be assumed to match its documented/kinematic specification
without direct measurement.

\section{Navigation Robustness in Tight Spaces}
\label{sec:results-cusp}

Table~\ref{tab:results-cusp-summary} restates
Section~\ref{sec:cusp-freeze}'s K-turn/reversal-cusp results: the
representative corridor reproduction goal's manoeuvre duration and
recovery count, across the untouched baseline and the final,
corner-extending \texttt{CuspStraightenerSmoother} configuration.

\begin{table}[htbp]
  \centering
  \caption{K-turn manoeuvre duration and recovery count, before and after the cusp-straightening smoother (simulation, three trials each).}
  \label{tab:results-cusp-summary}
  \begin{tabular}{@{}lcc@{}}
    \toprule
    & Baseline (no smoother) & Final config (smoother + corner extension) \\
    \midrule
    Mean duration & \SI{110.2}{\second} & \SI{43.9}{\second} \\
    Mean recoveries per trial & 2.0 & 0.0 \\
    Worst-case duration & \SI{168.1}{\second} & \SI{48.1}{\second} \\
    \bottomrule
  \end{tabular}
\end{table}

This is corroborated, not merely repeated, by real-hardware testing:
two independently sent goals on the physical vehicle, each requiring a
genuine reversal manoeuvre, both succeeded with zero recoveries
(Section~\ref{sec:cusp-hw-validation}) -- a stronger form of validation
than the simulation trial count alone, since it confirms the fix
transfers to real sensor noise, real actuator response, and real
planner behaviour on unseen goal geometry, not only the specific
corridor reproduction it was developed against. The one qualification
recorded against this result is that individual replans on real
hardware occasionally showed substantially higher cusp counts
(up to 8--10) than any simulation trial produced (maximum
$\approx$3) -- both hardware goals still succeeded despite this, so it
is reported as an open, unexplained gap between simulated and real
replanning behaviour rather than a limitation of the fix itself.

\section{Navigation Goal-Reaching}
\label{sec:results-goal-reaching}

Closed-loop \texttt{navigate\_to\_pose} goals requiring both straight
driving and real turning were completed successfully on the physical
QCar on multiple occasions across this project, confirmed both by
\texttt{bt\_navigator}'s own \texttt{Goal succeeded} status and, in
the final validation of Section~\ref{sec:final-validation}, by the
operator's direct observation that localisation tracked the vehicle's
real motion closely throughout the goal, rather than exhibiting the
map-shift artefact of Section~\ref{sec:display-lag}. The one
navigation timing measurement available (Table~\ref{tab:display-lag})
was taken under the pre-\texttt{readMode=0} bridge configuration and
should be read as an upper bound on the true post-fix pipeline
latency and stop-coast behaviour during a live goal, not as the final
figure, per the caveat in Section~\ref{sec:display-lag-interpretation}.

\section{Applicability of the Results}
\label{sec:results-applicability}

\subsection{Beyond This Specific QCar Unit}
\label{sec:applicability-unit}

The throttle deadband/gain values (Equation~\ref{eq:throttle-calibration})
and the steering-trim value (Section~\ref{sec:steering-trim}) are
per-unit calibration constants, expected to differ, potentially
substantially, on a different physical QCar chassis, or even on
this same chassis under different tyre wear, floor surface, or
battery charge conditions -- the calibration explicitly found this
floor to drift within a single session (Section~\ref{sec:deadband-discovery}).
The \emph{method} used to find them -- an empirical fixed-duty sweep
under real load, rather than an analytically derived or single-point
gain estimate -- is expected to transfer directly to any similarly
actuated small-scale ground vehicle, independent of the specific
numeric result.

\subsection{Beyond the QCar Platform}
\label{sec:applicability-platform}

The \texttt{readMode} buffering defect (Section~\ref{sec:readmode-bug})
is specific to the Quanser HAL's task-based I/O mode, but the general
failure class it represents -- a background acquisition task started
at object-construction time, with a read-consumption rate that
happens to match its production rate closely enough that a
start-up-gap backlog never drains -- is a generic ring-buffer hazard
that is not specific to this vendor's library, and is worth checking
for in any hardware abstraction layer offering a buffered or
task-based I/O mode as an alternative to immediate/polled reads,
particularly where the buffer is initialised eagerly rather than on
first read. The diagnostic principle that resolved it -- when an
internally self-consistent software abstraction (here, the
\texttt{Odometry} class) produces a value that contradicts direct
physical observation, verify the abstraction's own \emph{input}, not
only its internal consistency -- is a general debugging principle,
not specific to this project at all.

\subsection{Beyond Navigation}
\label{sec:applicability-domain}

The active-braking safety incident (Section~\ref{sec:active-braking})
and its resolution generalise directly to any reactive
feedback-control design intended for direct hardware deployment on a
system with confirmed actuator lag or inertia, regardless of
application domain: this project's standing rule -- trace new
reactive control logic through by hand for an overshoot scenario, or
validate it in simulation, before it reaches a physical actuator --
is a general control-engineering caution reinforced, not discovered,
by this project's own experience.
